\documentclass{article}

\usepackage{amsmath}
\usepackage{algorithmic}
\usepackage{graphicx}
\usepackage{xspace}

\begin{document}
\title{\Large Hello \LaTeX\ World}
\date{}
\author{Jucheol Moon\\\small Computer engineering and Computer Science\\\small California State University Long Beach\\\small\texttt{jucheol.moon@csulb.edu}}
\maketitle
\begin{abstract}
This document is a model and instructions for \LaTeX\ `article' class
\end{abstract}

\section{Introduction}
Welcome to the \LaTeX\ world.

\section{Ease of Use}

\subsection{Maintaining the Integrity of the Specifications}
The `article' class is used to format your paper and style the text. All margins, column widths, line spaces, and text fonts are prescribed.

\section{Styling Guide}

\subsection{Abbreviations and Acronyms}
Define abbreviations and acronyms the first time they are used in the text, 
even after they have been defined in the abstract.

\subsection{Equations}
\begin{equation}
\sum_{n=0}^\infty\frac{af^n}{n!}(x-a)^n
\end{equation}

\noindent(1) is the famous Taylor series. Use ``(1)", not ``Eq. (1)" or ``equation (1)", except at the beginning of a sentence: ``Equation (1) is . . ."
\newline\indent Taylor series in a text would be $\sum_{n=0}^\infty\frac{af^n}{n!}(x-a)^n$.

\subsection{Lists}
Bullet style list.
\begin{itemize}
\item item 1
\item item 2
\item item 3
\end{itemize}

Number style list.
\begin{enumerate}
\item item 1
\item item 2
\item item 3
\end{enumerate}

\subsection{Figures and Tables}
\paragraph{Positioning Figures and Tables} Figure captions should be below the figures; table heads should appear above the tables. Insert figures and tables after they are cited in the text. Use the abbreviation ``Fig. 1".
\begin{table}[!h]
\medskip
\begin{center}Table 1: Table Type Styles
\\\medskip\medskip
\begin{tabular}{|c||c|c|c|}
\hline \textbf{Table} & \multicolumn{3}{c|}{\textbf{Table Column Head}}\\
\cline{2-4} \textbf{Head} & \textbf{\textit{Table column subhead}} & \textbf{\textit{Subhead}} & \textbf{\textit{Subhead}}\\
\hline&&&\\
\hline
\end{tabular}
\end{center}
\end{table}

\begin{figure}[h!]
\centering
\includegraphics[width=0.4\textwidth]{fig1.png}
\caption{Working example}
\end{figure}

\subsection{Algorithms}
\begin{algorithmic}
\STATE $i\gets 10$
\IF{$i\geq 5$}\STATE $i\gets i-1$
\ELSE
\IF{$i\leq 3$}\STATE $i\gets i+2$
\ENDIF
\ENDIF
\end{algorithmic}

\subsection{Source codes}
\begin{verbatim}
public class HelloWorld {
   public static void main(String[] args) {
      System.out.println("Hello, World");
   }
}
\end{verbatim}

\subsection{References}
Please number citations consecutively within brackets [1]. The sentence punctuation follows the bracket [2]. Refer simply to the reference number, as in [3]---do not use ``Ref. [3]" or ``reference [3]" except at the beginning of a sentence.

\begin{thebibliography}{00}
\bibitem{Eason} G. Eason, B. Noble, and I. N. Sneddon, ``On certain integrals of Lipschitz-Hankel type involving products of Bessel functions,'' Phil. Trans. Roy. Soc. London, vol. A247, pp. 529--551, April 1955.
\bibitem{Maxwell} J. Clerk Maxwell, A Treatise on Electricity and Magnetism, 3rd ed., vol. 2. Oxford: Clarendon, 1892, pp.68--73.
\bibitem{Jacobs} I. S. Jacobs and C. P. Bean, ``Fine particles, thin films and exchange anisotropy,'' in Magnetism, vol. III, G. T. Rado and H. Suhl, Eds. New York: Academic, 1963, pp. 271--350
\end{thebibliography}

\end{document}
